\section{Conclusion and open problems}
\label{sec:conclusion}

The counting bound for complete caps has a slack, and the slack has a
geometry: every hyperplane carries the prescribed share
\(Q(s_\Pi)\) of it, so near the counting bound the hyperplane sections of
a complete cap have sizes in a set of two or three short runs.  The
classical character equations, over all hyperplanes and over the
hyperplanes through a secant, then become integer systems on a small
support, and their infeasibility, or the collision count they force on
every secant, excludes the cap.  In the coding dictionary the statement
is that a quasi-perfect \([k,k-d-1,4]_q\) code of length equal to the
sphere-covering bound has a weight distribution supported on a set that
the first two power moments cannot realize.  Five parameter sets in
dimensions four and six are excluded in this way, and the exact scan
shows that in the tested range the mechanism gains exactly one step.

Three questions remain.

\begin{enumerate}[label=(\arabic*),nosep]
\item \emph{Beyond one step.}  One step above the counting bound the
 admissible set is an interval and every test of this paper is
 satisfied.  The losses \(Q(s_\Pi)\) are then still confined to
 \([0,\Lambda_0]\), and their distribution over the hyperplanes through
 a secant, rather than the section sizes alone, is the information not
 yet used.  Is there a second-order constraint on that distribution that
 survives an interval of admissible sizes?
\item \emph{Higher moments.}  The third moment of the section sizes
 through a secant counts triples of further cap points by the dimension
 of their span with the secant, and brings in the five-subsets contained
 in a solid, the weight-five words of \(C_A^\perp\).  Together with a
 lower bound on that count it would sharpen Theorem~\ref{thm:local-bounds};
 no such bound is known to us.
\item \emph{The proportion problem.}  Proposition~\ref{prop:degree-countermodel}
 shows that the incidence identities alone cannot force collisions on a
 positive proportion of secants.  Can the hyperplane loss identity,
 perhaps in the form of Proposition~\ref{prop:holes} applied to
 sub-configurations, supply the missing constraint?
\end{enumerate}

The scan also raises a small empirical question: dimensions four and six
each yield two exclusions in the tested range and dimension five none.
The constant \(C_0=\binom{k-2}{2}\theta_{d-4}\) vanishes only for
\(d\le3\), so this is not an artefact of the moment formula; whether it
reflects the parity of \(d\) or only small numbers would be settled by
extending the scan in dimension five.
